PP is sensitive to perimeter measurement, which in turn depends on the
geographic projection used. This section documents the artifact and the
correction methodology.

\subsection{The TIGER/Line Boundary Artifact}

U.S.\ Census TIGER/Line shapefiles represent geographic boundaries at a
resolution of approximately 100--500 metres for tract boundaries, but coastlines
and rivers are digitised at higher detail, introducing small zigzag segments
that faithfully represent physical geography but inflate perimeter relative
to any smooth representation of the same boundary.

The consequence for PP is that coastal and riverine districts appear less compact
than their ``true'' geometric shape warrants.
Consider a district whose true boundary, if smoothed, would give PP $= 0.28$.
Digitised at TIGER/Line resolution, the same district may report PP $= 0.22$---a
21\% underestimate purely from boundary representation artifacts.

\subsection{Projection Effects}

Longitude-latitude coordinates (WGS84, EPSG:4326) distort area and distance
at mid-latitudes, making PP unreliable for direct comparison across states.
Computation in geographic coordinates (degrees) systematically biases perimeter
relative to area because one degree of longitude is shorter in absolute distance
than one degree of latitude at non-zero latitudes.

\paragraph{Recommended projection.}
All PP computations in the \textsc{bisect} platform use Albers Equal Area Conic
(EPSG:5070), the standard equal-area projection for the contiguous United States.
This projection:
\begin{itemize}
  \item Preserves area exactly (by definition of an equal-area projection).
  \item Minimises angular distortion at mid-latitudes ($\approx 23\degree$--$\approx 49\degree$N).
  \item Is the projection used by the U.S.\ Geological Survey for national maps
    and the standard for redistricting data in academic literature~\citep{census2020tiger}.
\end{itemize}

\subsection{Quantifying the Correction}

We compare PP computed in raw WGS84 and in EPSG:5070 for all NC, WI, and TX
districts under ratio-optimal structure (seed $= 42^\dagger$).
Table~\ref{tab:projection} reports the mean PP shift.

\begin{table}[ht]
\centering
\caption{Mean PP improvement from Albers Equal Area Conic (EPSG:5070) vs.\
raw WGS84 coordinates (ratio-optimal, seed $= 42^\dagger$).}
\label{tab:projection}
\begin{tabular}{lccc}
\toprule
State & Mean PP (WGS84) & Mean PP (EPSG:5070) & Shift \\
\midrule
NC ($k=14$) & 0.185 & 0.225$^\dagger$ & $+0.040$ \\
WI ($k=8$)  & 0.210 & 0.248$^\dagger$ & $+0.038$ \\
TX ($k=38$) & 0.168 & 0.203$^\dagger$ & $+0.035$ \\
\bottomrule
\end{tabular}
\end{table}

The projection correction is largest for NC because several NC congressional districts
border the Atlantic coast and the Outer Banks, where TIGER/Line digitisation of
the shoreline is most detailed.
Coastal states such as FL, LA, and MA would exhibit larger corrections;
inland states such as KS and NE would exhibit corrections near zero.

\subsection{Boundary Smoothing (Not Recommended)}

An alternative to projection correction is boundary smoothing---applying a
Douglas-Peucker algorithm~\citep{douglas1973algorithms} to reduce vertex count.
We do \emph{not} recommend this approach for legal applications because:
(1) smoothing is not standardised and opposing counsel can challenge the smoothing
parameter as arbitrary; (2) smoothing removes geometric information and may obscure
genuine non-compactness in highly irregular districts; and (3) projection to
EPSG:5070 achieves the necessary correction without altering the underlying geometry.
